\section{Primitive convention and character-separated leakage}
\label{sec:primitive}

Primitive closed words are quotiented by cyclic rotation, never by
reflection.  For a primitive class $[\gamma]$ set
\[
 w(\gamma)=\prod_j w(S_{j+1}),\qquad
 H_\alpha(\gamma)=\prod_j\alpha(S_j,S_{j+1}).
\]
Formally,
\[
 D_{\rho,P}
 =\prod_{[\gamma]\ \mathrm{primitive}}
   \det\!\left(I-w(\gamma)\rho(H_\alpha(\gamma))\right).
 \tag{5.1}\label{5.1}
\]
The $m$-fold traversal uses both $w(\gamma)^m$ and
$H_\alpha(\gamma)^m$.  Koszul signs stay ordinary scalar coefficients; no
supertrace is introduced.

An unmarked monomial such as $x^3y^3$ generally aggregates many primitive
words and temporal repetitions.  To isolate a particular cycle, attach a
commuting marker $z_{S,T}$ to each directed edge.  This works when the edge
multiset has a unique connected cyclic traversal; otherwise a finer
cyclic-word marker is required.

\begin{theorem}[Separated holonomy forces a character leak]
\label{thm:separated-leak}
Let $\gamma$ be a primitive directed cycle isolated by such a marker.  If
$H_\alpha(\gamma)$ is not conjugate to the reference holonomy
$H_0(\gamma)$, then some irreducible block has a different first-traversal
coefficient.  If $H_0(\gamma)=e$ and $H_\alpha(\gamma)\ne e$, then
\[
 \chi_\rho(H_\alpha(\gamma))\ne d_\rho
 \tag{5.2}\label{5.2}
\]
for at least one irreducible $\rho$.
\end{theorem}

Indeed, the trace-log identity
\[
 \log\det(I-B_\rho)
 =-\sum_{n\ge1}\frac{\tr(B_\rho^n)}n
 \tag{5.3}\label{5.3}
\]
assigns the isolated first traversal the coefficient
$-\chi_\rho(H_\alpha(\gamma))w(\gamma)z_\gamma$.  Irreducible characters
separate conjugacy classes.  The full proof, including the identity-reference
case, is in \cref{app:leak-proof}.

\Cref{thm:separated-leak} is deliberately one-way.  It does not assert that
unmarked character determinants classify cocycles, and it does not prevent
collisions among aggregated orbit data.
